Language-model agents increasingly negotiate and transact with one another, yet evaluations usually score a final outcome or an isolated decision. We study twenty complete four-player Monopoly games played by four configured LLM systems, linking every offer, construction attempt, financing route and bankruptcy to the engine's legal menus and recorded events. Recipients resist completing a rival's color group. Offers that would complete the offer author's group are accepted 19 times in 289 and rejected 159 times, although the median such offer pays the recipient \$200. Counters to these offers rarely remove the completing deed: 92 of 111 keep it, and nearly all that give the counterer nothing ask for more cash. Three configurations accept such offers less often than offers completing their own group; the fourth, which ran mostly under a refusal-recovery condition, rarely rejects but prices the completion in its counters. Completion becomes construction in 43 of 63 trade-completed groups, and in one configuration most of the shortfall comes from engine defaults on builds that were executable after legal reordering. Nominal wealth misstates liquidity in both directions: three eliminated players could legally have paid the debt that ended their game, while one with \$1,310 of nominal assets could reach \$440. Creditor inheritance sends 54 of 60 estates to the eventual winner, and the most consequential bargaining conduct in the selected transcripts occurs where solvency and inheritance are at stake. The study is observational and estimates no causal effects, optimal play or model ranking.
